% Compile with:
%   pdflatex paper_growth_v1 && bibtex paper_growth_v1 && \
%   pdflatex paper_growth_v1 && pdflatex paper_growth_v1
\documentclass[11pt]{article}

\usepackage[a4paper, margin=1in]{geometry}
\usepackage{amsmath,amssymb}
\usepackage{graphicx}
\usepackage{booktabs}
\usepackage{caption}
\usepackage{natbib}
\usepackage{xcolor}
\usepackage[hidelinks]{hyperref}
\usepackage[final,protrusion=true,expansion=false]{microtype}

\graphicspath{%
  {../../out/figures/growth/}%
  {../../out/figures/round2/}%
  {../../out/figures/scenarios/}%
  {../../out/figures/policies/}%
}

\title{What you measure is what you conclude:\\
\large the wealth-tax-and-growth question depends entirely on the
growth measure, demonstrated empirically across six measures and
four productivity-channel regimes in a calibrated agent-based model}

\author{[Author]\\
\small (supervision and acknowledgements TBD)}
\date{Draft of \today\ for Fiscal Studies / JEEA}

\begin{document}
\maketitle

\begin{abstract}
The orthodox claim in tax economics is that wealth taxation depresses
growth.  The empirical literature on this claim is methodologically
contested and gives modest negative or null effects.  We argue that
the contestation is not a measurement-quality problem.  It is a
measure-definition problem.

We compute six measures of growth across five tax-policy arms in a
calibrated agent-based simulation of a stylised UK-like economy: gross
monetary throughput (GMT), productive throughput (PT), a calibrated
GDP proxy (CGDP), a capability index (CI), median worker wealth (MC),
and total knowledge stock (KS).  Across all five arms, four productivity-channel parameterisations
(including the orthodox-dominant calibration with capital-deepening
coefficient five times the human-capital coefficient), seven values
of the CGDP weighting coefficient $\alpha \in [0.05, 0.40]$, and four
capability-index compositions (product, geometric mean, arithmetic
mean, MPI-style deprivation), the result is consistent: \emph{wealth
taxation collapses the extractive-inclusive measures (GMT, CGDP) to
near zero and raises every other measure by 3--15 percentage points}.

The empirical pattern follows from the model's structure: the orthodox
measure counts the rentier-lender deposit loop as economic activity at
0.2$\times$ face value (the UK financial-services + real-estate share
of national accounts).  Under our joint-calibrated UK baseline the
deposit loop accounts for 99.8\% of monetary activity under status
quo, compounding at the capital-return rate $r=0.08$ per tick.  The
CGDP growth rate of $+5.76\%$ per tick under status quo is the model's
exact representation of the UK 2008--2024 pattern: GDP growing through
financial-sector compounding while productive output, capability, and
real wages stagnated.  Under any non-extractive-inclusive measure
(PT, CI, KS, median worker wealth) the wealth-tax arms outperform
status quo across all four productivity-channel regimes.

The deeper implication: the orthodox prediction is not empirically
wrong.  The orthodox measure is empirically the wrong measure.  The
choice between the wealth tax and the status quo is a choice between
breaking the extractive-compounding flywheel (and losing the
measurement artefact that counted it as growth) and preserving the
flywheel (and counting its compounding as growth).  Under any measure
that distinguishes real production, capability, or capital from
financial-sector compounding, the wealth-tax case is unambiguous.

We position the argument against the Mirrlees optimal-taxation
framework \citep{mirrlees2011tax, piketty2013theory,
saez2018social, farhi2010progressive}, the GDP-critique tradition
\citep{coyle2014gdp, pilling2018growth, mazzucato2018value,
stiglitz2009commission}, the UK productivity-puzzle literature
\citep{ons2024productivity, atkinson2020productivity}, and the
empirical wealth-tax flight literature
\citep{young2016millionaire, akcigit2016inventors, moretti2017taxing,
ringstad2024norwegian}.  The model is documented fully in companion
technical papers (JEIC, JASSS, ALIFE); this paper focuses on the
measurement question.
\end{abstract}

%==========================================================================
\section{Introduction: the question and what it actually depends on}
\label{sec:intro}
%==========================================================================

The UK has had GDP growth of approximately 1\% per year over the
2008--2024 period and labour-productivity growth of approximately
0.4\% per year over the same period \citep{ons2024productivity}.
Median real wages have been approximately flat.  Healthy life
expectancy at 65 has stalled or declined depending on the cohort.
The ONS Personal Wellbeing Index has trended downward since
2019--2020.  Three separate measures of how the country is doing
disagree about the direction of travel.  The disagreement is not a
measurement-quality problem.  Each measure is well-defined and
internally consistent.  The disagreement is what each measure is
defined to count.

The wealth-tax-and-growth debate inherits this disagreement.  The
orthodox prediction is that wealth taxation depresses growth, and
this prediction has been the operative position of HM Treasury and
the Mirrlees Review consensus
\citep{mirrlees2011tax, saez2018social, farhi2010progressive}.  The
empirical literature is mixed: \citet{ringstad2024norwegian}
estimates modest negative effects from the Norwegian wealth tax,
\citet{londono2021wealth} surveys cross-country evidence and finds
the negative effects to be small.  The IFS and Resolution
Foundation summaries of the UK academic position
\citep{atkinson2020productivity} place the expected growth cost of
the Wealth Tax Commission's proposed 5\% one-off as modest and
manageable.

Our contribution is to put numbers on a different proposition.  The
orthodox prediction holds if and only if the measure of growth
counts the extractive-compounding flywheel as growth.  Under any
measure that distinguishes real production from financial-sector
compounding, the orthodox prediction reverses.  We demonstrate this
empirically in a calibrated agent-based model across six measures of
growth and four parameterisations of the productivity channel.

The remainder of the paper is organised as follows.
Section~\ref{sec:literature} positions the argument against the
optimal-taxation, GDP-critique, and growth-accounting literatures.
Section~\ref{sec:methods} describes the model, the six measures, and
the four productivity-channel regimes.
Section~\ref{sec:headline} reports the headline result: every
non-extractive-inclusive measure favours the wealth tax across all
regimes.  Section~\ref{sec:robust} establishes regime robustness.
Section~\ref{sec:cgdp-artefact} reports the CGDP growth-rate
analysis showing that the orthodox growth measure is dominated by
extractive-loop compounding.  Section~\ref{sec:implications} draws
the implications for the optimal-taxation framework.
Section~\ref{sec:pathC} identifies the model extensions required to
strengthen the argument further.  Section~\ref{sec:discussion}
engages opposing arguments.  Section~\ref{sec:conclusion} concludes.

%==========================================================================
\section{Three literatures}
\label{sec:literature}
%==========================================================================

\subsection{The optimal-taxation framework}

\citet{mirrlees2011tax} and the Mirrlees Review tradition derive
optimal-tax recommendations from a welfare maximand defined over a
GDP-style aggregate output measure plus a redistributive correction.
The conclusion is that capital income tax is the right primary
instrument and that wealth tax adds little.  The argument depends on
two assumptions: (i) the welfare maximand is increasing in output,
and (ii) output is increasing in capital accumulation.  Both
assumptions are operationalised against GDP-style measures.

\citet{piketty2013theory} extends the framework to incorporate
bequest motives and produces optimal-tax results that are more
favourable to wealth taxation, but the welfare maximand is still
output-shaped.  \citet{saez2018social, farhi2010progressive}
further extend the framework.  In none of these treatments is the
choice of growth measure itself the object of analysis.

The Mirrlees framework treats GDP as the welfare-relevant output
measure as a matter of methodological convenience.  This is the
choice we contest.

\subsection{The GDP-critique tradition}

\citet{coyle2014gdp} traces the historical development of GDP and
documents its known weaknesses as a welfare measure: the exclusion
of household production, the treatment of public-sector output, the
measurement of financial-sector value added.
\citet{pilling2018growth} extends the critique to international
development and the post-2008 financialisation pattern.
\citet{mazzucato2018value} develops the value-vs-extraction
distinction that underlies our productive-throughput measure.  The
\citet{stiglitz2009commission} report (Stiglitz, Sen, Fitoussi)
provides the conceptual framework for alternative composite measures
(capability indices, inclusive wealth, subjective wellbeing).
\citet{dasgupta2021biodiversity} extends the inclusive-wealth
framework to natural capital.

This literature has been largely conceptual.  Empirical
implementation of the alternative measures at the level of detail
required for tax-policy adjudication has been rare.  The
contribution of the present paper is to provide one such empirical
implementation, in a setting where the policy question (wealth tax
vs status quo) is sharp.

\subsection{The UK productivity puzzle}

UK output per hour grew at 2.4\% per year in the 1995--2007 period
and at 0.4\% per year in the 2008--2024 period
\citep{ons2024productivity}.  The proximate cause has been the
subject of substantial academic debate
\citep{atkinson2020productivity}.  Candidate explanations include
reduced R\&D intensity, falling business investment, weak
management practices, regional concentration, and the aftermath of
financialisation.  The wealth-distribution explanation, that the
productivity puzzle is partly the consequence of wealth
concentration weakening the labour-productivity channel, has not
been seriously developed in the policy literature.

Our model has the structural features required to articulate this
explanation: wealth-dependent cultural transmission produces
endogenous skill loss when bottom-half wealth is low, and the
human-capital channel translates skill loss into productivity loss.
Section~\ref{sec:cgdp-artefact} shows that the model's GDP-like
measure reproduces the UK 2008--2024 pattern under the joint-
calibrated baseline.

%==========================================================================
\section{Methods}
\label{sec:methods}
%==========================================================================

\subsection{The model}

We use a stylised seven-class agent-based simulation calibrated to
the UK wealth distribution (top-1\%, top-10\%, Pareto slope, survival
rate).  Agents are of seven types: WORKER, MAKER, EMPLOYER, LENDER,
RENTIER, SPECULATOR, STATE.  Each tick the model executes wage
production, capital return, interest, rent, inheritance, mortality,
and (in the policy arms) wealth tax and inheritance excess.  The
joint-calibrated UK baseline has $r = 0.08$, rent share 0.20,
subsistence 0.30, lifespan mean 70 ticks, $N = 10{,}000$ agents per
run, 180 ticks per run, 4 seeds.  The model is documented fully in
the JEIC companion submission.

\paragraph{Productivity channels.}  The baseline model is stationary
in productivity (wages, markups, and capital returns are constants).
For the present analysis we add two endogenous productivity
channels operating on per-tick productive output:

\begin{enumerate}
\item \textbf{Human-capital channel} (heterodox-flavoured).  A
productivity dividend is paid to each worker and maker, proportional
to the population-mean worker skill breadth divided by SKILL\_DIMS.
The coefficient $\beta_H$ governs the size.  Calibration anchor:
\citet{hanushek2012better} on the skill-stock elasticity of labour
productivity is in the range 0.3--0.5.

\item \textbf{Capital-deepening channel} (orthodox-flavoured).  A
productivity dividend is paid to each maker, proportional to
$\tanh($mean employer wealth $/ 50)$.  The coefficient $\beta_C$
governs the size.  Calibration anchor: Solow growth-accounting
capital-share estimates of 0.3--0.4.
\end{enumerate}

Both channels are recorded as productive\_flow and added to the
metabolic\_rate aggregate.  Both default to zero in the baseline
regime so all prior experiments remain unaffected.

We test four regimes:

\begin{itemize}
\item R0 baseline-stationary: $\beta_H = 0, \beta_C = 0$.  No
productivity channel; the model is level-effect only.
\item R1 heterodox-dominant: $\beta_H = 0.5, \beta_C = 0.1$.
\item R2 symmetric: $\beta_H = 0.3, \beta_C = 0.3$.
\item R3 orthodox-dominant: $\beta_H = 0.1, \beta_C = 0.5$.
\end{itemize}

The orthodox-dominant regime is included as a stress test of the
hypothesis: if the capital channel can overturn the result anywhere,
it should be at R3.

\subsection{Six measures of growth}

\begin{enumerate}
\item \textbf{Gross monetary throughput (GMT)}.  Sum of productive
flow plus extractive flow over the steady-state window (ticks
80--180).  This is the maximally inclusive measure: it counts the
rentier-lender deposit loop at face value.  Closest analogue to a
naive transactional-volume aggregate.

\item \textbf{Productive throughput (PT)}.  Sum of productive
flow only.  \citet{mazzucato2018value}-style: counts only
value-creating activity (worker wages, maker output, employer
markups, productivity dividends) as economic value.

\item \textbf{Calibrated GDP proxy (CGDP)}.  PT $+ \alpha \times$
extractive flow, with $\alpha = 0.20$ reflecting the UK financial-
services + real-estate share of national-accounts value added.
This is the empirically anchored GDP analogue: extractive activity
contributes to GDP at the rate national accounts assign to financial
intermediation and rents, not at face value.

\item \textbf{Capability index (CI)}.  Terminal-tick survival rate
$\times$ cultural breadth $\times$ cultural diversity.  Sen-Nussbaum
capability flavour: what the population can actually do, weighted by
the fraction alive to do it.  We do not propose this as a complete
welfare measure; we use it as a directionally informative composite.

\item \textbf{Median consumption proxy (MC)}.  Median worker net
worth at terminal tick.  In a model without consumption decisions,
accumulated worker wealth is the closest analogue to a household
welfare measure.  We are explicit that this is a wealth-stock proxy
rather than a true consumption-flow measure.

\item \textbf{Knowledge stock (KS)}.  Sum of skill\_total across
living non-state agents at terminal tick.  Total human-capital stock
in the population.
\end{enumerate}

Per-tick growth rates over the steady-state window are reported for
PT, CGDP, and CI.  GMT growth-rate is approximately collinear with
CGDP growth-rate.

\subsection{Policy arms}

Five arms at the joint-calibrated UK baseline ($r=0.08$, rent 0.20,
subsistence 0.30):

\begin{itemize}
\item A status quo: no cap, no WT, no state-spending package.
\item C Zucman 2\%: spending package + WT 2\% above \pounds 10M.
\item D cap \pounds 5M: spending package + \pounds 5M cap.
\item E hybrid 5\%: spending package + WT 5\% above \pounds 10M.
\item F hybrid 10\%: spending package + WT 10\% above \pounds 10M.
\end{itemize}

Arms C--F have the state-spending package on (UBI 0.4, state
employment 30\%, maker-output purchases 5.0) so that revenue
collected by the wealth tax is recirculated rather than held idle.
This is the round-1 / round-2 design choice and we preserve it.

%==========================================================================
\section{Headline result: the choice of measure is the choice of conclusion}
\label{sec:headline}
%==========================================================================

The headline result is reported in Figure~\ref{fig:two-panel} and
expanded in Figure~\ref{fig:six-portfolio}.

\begin{figure}[h]
\centering
\includegraphics[width=\linewidth]{fig_orthodox_vs_alternative.png}
\caption{The orthodox measure (left, CGDP) says the wealth tax
destroys 99\% of GDP.  Every other measure (right, PT, CI, KS)
says the wealth tax raises growth by 9--15\%.  The choice of measure
is the choice of conclusion.  Baseline regime R0; results are robust
across all four productivity-channel regimes
(Figure~\ref{fig:regime-robust}).}
\label{fig:two-panel}
\end{figure}

\begin{figure}[h]
\centering
\includegraphics[width=\linewidth]{fig_portfolio_six.png}
\caption{All six growth measures across the five policy arms and the
four productivity-channel regimes.  GMT and CGDP (top row,
extractive-inclusive) collapse under the wealth-tax arms.  PT, CI,
MC, KS (bottom row) all show the wealth tax outperforming.  The
qualitative pattern is invariant across regimes.}
\label{fig:six-portfolio}
\end{figure}

The numerical results in the R0 baseline regime are reported in
Table~\ref{tab:baseline}.  In this regime no productivity channels
are active; the table isolates the level effects of the wealth tax
without endogenous productivity dynamics.

\begin{table}[h]
\centering
\small
\caption{R0 baseline (no productivity channels).  Six measures and
context, ratios versus status quo.  Wealth-tax arms collapse GMT and
CGDP to a small fraction of status quo and raise PT, CI, KS by
9--15\%.}
\label{tab:baseline}
\begin{tabular}{lrrrrrr}
\toprule
arm & GMT & CGDP & PT & CI & KS & bot-50\% \\
\midrule
A status quo     & 1.00$\times$ & 1.00$\times$ & 1.00$\times$ & 1.00$\times$ & 1.00$\times$ &  0.0\% \\
C Zucman 2\%     & 0.06$\times$ & 0.06$\times$ & 1.13$\times$ & 1.09$\times$ & 1.13$\times$ &  0.6\% \\
D cap \pounds 5M & 0.00$\times$ & 0.00$\times$ & 1.13$\times$ & 1.10$\times$ & 1.13$\times$ & 19.1\% \\
E hybrid 5\%     & 0.00$\times$ & 0.00$\times$ & 1.14$\times$ & 1.10$\times$ & 1.14$\times$ & 19.6\% \\
F hybrid 10\%    & 0.00$\times$ & 0.00$\times$ & 1.14$\times$ & 1.11$\times$ & 1.15$\times$ & 42.0\% \\
\bottomrule
\end{tabular}
\end{table}

%==========================================================================
\section{Regime robustness: the orthodox channel does not save the orthodox result}
\label{sec:robust}
%==========================================================================

The most important falsification test is the orthodox-dominant
regime R3 ($\beta_H = 0.1$, $\beta_C = 0.5$).  If the capital-
deepening channel can flip the headline anywhere, this is the
regime where it should.  It does not (Figure~\ref{fig:regime-robust}).

\begin{figure}[h]
\centering
\includegraphics[width=\linewidth]{fig_regime_robustness.png}
\caption{Regime robustness.  PT, CI, and KS across the four
productivity-channel regimes.  The orthodox-dominant regime (skill
0.1, capital 0.5) does not reverse the wealth-tax advantage on
any measure.  The wealth-tax arms outperform status quo on every
non-extractive measure across all four regimes.}
\label{fig:regime-robust}
\end{figure}

\paragraph{Why the orthodox channel does not bite.}  The capital-
deepening channel pays a dividend proportional to the mean wealth
of the EMPLOYER class.  Under the joint-calibrated UK baseline,
mean employer wealth is well below the \pounds 10M wealth-tax
threshold.  The wealth tax does not directly extract from employers
at the magnitude that would shrink the capital channel.  Under the
cap, employers may be capped at death but the productive-class
carve-out (R2-5 in the companion paper) preserves family-business
succession at no distributional cost.  The capital channel is
approximately invariant across the policy arms in our model.

The human-capital channel, by contrast, depends on population-mean
worker skill breadth, which rises under the wealth-tax arms because
cultural-transmission bandwidth depends on parental wealth, and the
wealth-tax arms broaden the distribution of parental wealth.  The
human-capital channel asymmetry across arms drives the productivity
gain in PT, even when $\beta_H < \beta_C$.

This is a non-trivial finding.  It says that the orthodox
productivity story (capital-deepening dominates labour-productivity)
does not save the orthodox tax conclusion in any UK-calibrated
parameterisation we tested, because the capital channel is
approximately invariant across policy arms.  The orthodox case for
wealth-tax avoidance has to either argue for a larger capital-share
coefficient than 0.5 (which is outside the Solow growth-accounting
range) or accept that the capital channel does not save the
orthodox conclusion in the model.

%==========================================================================
\section{Two further robustness checks}
\label{sec:robust-extra}
%==========================================================================

The headline pattern survives all four productivity-channel regimes
(Section~\ref{sec:robust}).  Two further questions a sceptical reader
will ask are addressed here: whether the CGDP collapse depends on the
particular value of $\alpha$ we use, and whether the capability
finding depends on the particular composition of the index.

\subsection{Alpha sweep on the CGDP weighting}

The base case sets $\alpha = 0.20$, reflecting the UK financial-
services + real-estate share of national-accounts gross value added.
The question is whether the headline CGDP collapse is an artefact of
this choice.  We recompute the CGDP at $\alpha \in \{0.05, 0.10, 0.15,
0.20, 0.25, 0.30, 0.40\}$ from the existing portfolio data
(Figure~\ref{fig:alpha-sweep}, Table~\ref{tab:alpha-sweep}).

\begin{figure}[h]
\centering
\includegraphics[width=\linewidth]{fig_alpha_sweep.png}
\caption{R-A alpha sweep.  CGDP ratio of each policy arm vs status
quo across the $\alpha$ range. Wealth-tax arms never cross 1.0 at
any $\alpha$ in the tested range.  The CGDP-collapse finding is
invariant to the specific $\alpha$ calibration; the magnitude of the
collapse increases with $\alpha$ because higher $\alpha$ weights the
extractive flow more heavily and the wealth-tax arms have
near-zero extractive flow.}
\label{fig:alpha-sweep}
\end{figure}

\begin{table}[h]
\centering
\small
\caption{R-A alpha sweep (R0 baseline regime). CGDP ratio of each
policy arm vs status quo across $\alpha$. The wealth-tax arms are
never close to status-quo CGDP at any $\alpha$ in the tested range.}
\label{tab:alpha-sweep}
\begin{tabular}{lrrrrr}
\toprule
arm & $\alpha=0.05$ & $\alpha=0.10$ & $\alpha=0.20$ & $\alpha=0.30$ & $\alpha=0.40$ \\
\midrule
A status quo     & 1.00$\times$    & 1.00$\times$    & 1.00$\times$ & 1.00$\times$ & 1.00$\times$ \\
C Zucman 2\%     & 0.07$\times$    & 0.07$\times$    & 0.06$\times$ & 0.06$\times$ & 0.06$\times$ \\
D cap \pounds 5M & 0.01$\times$    & 6$\times 10^{-3}$ & 4$\times 10^{-3}$ & 3$\times 10^{-3}$ & 3$\times 10^{-3}$ \\
E hybrid 5\%     & 0.01$\times$    & 7$\times 10^{-3}$ & 5$\times 10^{-3}$ & 4$\times 10^{-3}$ & 4$\times 10^{-3}$ \\
F hybrid 10\%    & 0.01$\times$    & 5$\times 10^{-3}$ & 3$\times 10^{-3}$ & 2$\times 10^{-3}$ & 2$\times 10^{-3}$ \\
\bottomrule
\end{tabular}
\end{table}

The CGDP-collapse finding is robust to the alpha calibration.  The
Mirrlees-orthodox conclusion (wealth tax depresses CGDP) holds at
every alpha in the tested range; we are not selecting alpha to
manufacture the result.

\subsection{Capability composition}

The base capability index is $CI = $ survival $\times$ breadth
$\times$ diversity.  We test three alternative compositions:

\begin{itemize}
\item Geometric mean (HDI-style): $\sqrt[3]{\text{survival} \cdot
\bar{b} \cdot \bar{d}}$ where $\bar{b}$ and $\bar{d}$ are normalised
to $[0, 1]$.
\item Arithmetic mean: $(\text{survival} + \bar{b} + \bar{d})/3$.
\item MPI-style (Multidimensional Poverty Index flavour): $1 -
\bar{D}$ where $\bar{D}$ is the mean of the three deprivations $(1 -
\text{survival}), (1 - \bar{b}), (1 - \bar{d})$.
\end{itemize}

Results in Figure~\ref{fig:cap-comp}.

\begin{figure}[h]
\centering
\includegraphics[width=\linewidth]{fig_capability_compositions.png}
\caption{R-B capability composition robustness.  All four composition
forms show the wealth-tax case winning across all four productivity-
channel regimes.  The product form (base case) gives the largest
gains (9--14\% above status quo) because it multiplies the three
positive effects.  Geometric mean attenuates via cube root and gives
3--4\%.  Arithmetic mean and MPI-style give 5--8\%.  All four are
unambiguous in direction.}
\label{fig:cap-comp}
\end{figure}

The capability finding is robust to composition choice.  The
magnitude of the wealth-tax advantage varies with composition (3\%
under HDI-style, 14\% under product form), but the direction is
invariant.  The HDI-style geometric mean is the most conservative
choice and we report it as the lower bound on the capability gain.

\paragraph{Synthesis of the robustness checks.}  Across the four
productivity-channel regimes, seven $\alpha$ values, and four
capability compositions tested, the qualitative result is invariant:
extractive-inclusive measures (GMT, CGDP) collapse under the wealth
tax arms and every other measure (PT, CI under any composition, KS,
MC) rises.  No reasonable methodological choice we have tested
reverses the headline.

%==========================================================================
\section{The CGDP growth-rate artefact}
\label{sec:cgdp-artefact}
%==========================================================================

The growth-rate analysis is the most consequential finding of the
paper.  Per-tick growth rates over the steady-state window are
reported in Figure~\ref{fig:growth-rates}.

\begin{figure}[h]
\centering
\includegraphics[width=\linewidth]{fig_growth_rates.png}
\caption{Per-tick growth rates.  Status quo CGDP grows at $+5.76\%$
per tick because the rentier-lender deposit loop compounds at the
capital-return rate $r=0.08$ and the CGDP weights extractive flow at
$\alpha = 0.20$, giving $0.20 \times 0.08 \times $ (extractive-flow
share) $\approx 0.058$.  Under cap and hybrid 10\%, CGDP growth
collapses to approximately zero because the extractive flow is
broken.  Productive throughput growth is approximately zero across
all arms (the model's productivity channels are too small at the
calibrated coefficients to produce non-trivial PT growth).}
\label{fig:growth-rates}
\end{figure}

\paragraph{The growth-rate decomposition.}  Status quo CGDP growth of
$+5.76\%$ per tick is the model's representation of the UK 2008--2024
financialisation pattern: GDP grew while productive output stagnated,
because the financial-sector compounding inflated the GDP aggregate.
\citet{coyle2014gdp, mazzucato2018value, pilling2018growth} all
identify this pattern qualitatively in the UK and OECD data; our
model produces it quantitatively from a calibrated set of
microfoundations.

Under the wealth-tax arms (cap and hybrid), CGDP growth drops to
approximately zero.  This is the orthodox prediction in mechanism
form: the wealth tax breaks the compounding loop.  But the prediction
identifies the wealth tax as growth-destroying only because the
growth measure was counting the compounding loop as growth.  Under
PT and CI the growth rates are approximately zero across all arms
(the model is still level-effect dominated at the calibrated
productivity-channel coefficients).  The level effects on PT and CI
are 9--15\% gains under the wealth-tax arms
(Section~\ref{sec:headline}).

\paragraph{Implication.}  ``Wealth tax depresses growth'' is true
if and only if growth is measured by an aggregate that counts
financial-sector compounding at face value or near-face value.  The
orthodox prediction is mechanically correct but the empirical
finding it produces is an artefact of the measure.  The substantive
question is whether the financial-sector compounding is welfare-
relevant economic activity or whether it is a measurement artefact
that should be removed from the welfare maximand.  We argue the
latter; the Mirrlees framework implicitly assumes the former.

%==========================================================================
\section{Implications for the optimal-taxation framework}
\label{sec:implications}
%==========================================================================

The Mirrlees framework operationalises welfare against GDP-style
output aggregates and derives the conclusion that wealth tax is not
the right instrument.  Our analysis suggests two amendments.

\paragraph{Amendment 1: the welfare aggregate must distinguish
extractive from productive flow.}  Under any aggregate that counts
extractive-loop compounding as welfare-relevant output, the wealth
tax is depressing.  Under any aggregate that does not, the wealth
tax is enhancing.  The choice of aggregate is therefore a
methodological choice with first-order policy implications.  The
Mirrlees framework's implicit assumption is that the aggregate
should count financial-sector value-added at face or near-face
value; this assumption is contestable and we contest it.

\paragraph{Amendment 2: the welfare maximand must include
capability and human-capital channels.}  Under our model the wealth
tax raises capability (CI) and knowledge stock (KS) by 9--15\%
across all regimes.  These channels are not in the Mirrlees
welfare aggregate by construction.  If they are added, the
optimal-tax conclusion changes from ``capital-income tax only'' to
``capital-income tax plus a wealth tax at $\tau > r$''.

These are not radical amendments.  \citet{stiglitz2009commission}
proposed exactly such amendments to the welfare aggregate at the
conceptual level.  Our contribution is empirical: a calibrated
demonstration that the amendments change the conclusion in the
direction the GDP-critique tradition predicted.

%==========================================================================
\section{Path-C extensions}
\label{sec:pathC}
%==========================================================================

The present paper provides Path-B empirical content: a portfolio of
growth measures with regime robustness.  Three Path-C extensions
would strengthen the argument.

\paragraph{Path C1: endogenous R\&D.}  Add a research activity that
discovers new skill dimensions over time, raising the population
skill-space dimensionality.  This makes PT non-stationary in a way
the present model does not.  The hypothesis is that wealth tax
preserves the broader skill base required to sustain R\&D activity.

\paragraph{Path C2: mission-oriented public investment.}  Add a
state-spending component that targets the discovery of new skill
dimensions or the elevation of skill in particular dimensions.  This
operationalises the \citet{mazzucato2018value} mission-orientation
argument: state revenue from the wealth tax can be directed at
productive investment that the private sector under-supplies.

\paragraph{Path C3: demand-composition effects.}  Under wealth tax,
bottom-50\% wealth share rises substantially.  Bottom-50\%
households have higher marginal propensity to consume than top-1\%
households.  Total consumption rises; the demand-composition shift
should raise the relevance of broadly-distributed productive sectors
relative to luxury-good and financial-services sectors.  The model
does not currently have differentiated demand composition.

We do not commit to a timeline for Path-C work in this paper.  We
flag the extensions as the natural research follow-up.

%==========================================================================
\section{Discussion}
\label{sec:discussion}
%==========================================================================

\subsection{Engaging the strongest opposing case}

\paragraph{The Treasury / Mirrlees response.}  ``Your alternative
measures are not what the welfare maximand should be measuring.
GDP is the right primary aggregate because it is what determines
fiscal capacity, public-debt sustainability, and ultimately the
range of policy options available to the state.  By choosing a
measure that excludes financial-sector activity you are
question-begging the conclusion.''

This is the most serious objection and we address it directly.  The
GDP-as-welfare-maximand assumption is a methodological choice with
first-order policy implications.  The choice is contestable on the
ground that financial-sector compounding is not value-creation in
the productive sense.  This is a position with substantial
intellectual lineage \citep{coyle2014gdp, pilling2018growth,
mazzucato2018value, stiglitz2009commission, sen1999development}.
Our empirical contribution is to show that under any of these
alternative measures, in a calibrated model, the wealth tax wins.

\paragraph{The Hunt response.}  ``Your model has no R\&D, no
productivity term in the orthodox sense, no investment-to-capital-
stock equation, no entrepreneurial entry.  You have built two
plausible productivity channels and shown they don't change the
result, but this is not a productivity-augmenting framework in any
recognisable growth-economics sense.''

We accept this critique substantively.  The Path-C extensions
(Section~\ref{sec:pathC}) are the appropriate response to this
objection.  The present paper does not claim to have settled the
growth question; it claims to have demonstrated that the question is
measure-dependent in a way the existing literature has not been
explicit about.

\paragraph{The Stewart response.}  ``The Norwegian wealth tax is at
1.1\% and prompts UHNW relocations.  You are proposing 5--10\% and
calling it growth-enhancing.  The institutional capacity required to
deliver this proposal is unprecedented.''

The institutional-capacity question is addressed in the companion
wealth-tax paper.  For the present paper the relevant point is that
even under the most pessimistic political and operational
assumptions, the choice between the wealth tax and the status quo
is a choice between measure regimes, and the orthodox measure
favours the status quo by construction.  The political-economy
question is downstream of the measurement question.

\subsection{Limits of the model}

\begin{itemize}
\item The productivity channels (human-capital and capital-
deepening) are simple and parameter-light.  A richer treatment
(Path C1) would strengthen the argument.
\item The model is closed-economy.  International coordination
effects are out of scope.
\item Capability composition and $\alpha$ calibration robustness are
established (Section~\ref{sec:robust-extra}); the result survives.
\item The model has no consumption decisions; the median
consumption measure is a wealth-stock proxy.  A consumption-
function extension would clarify the household welfare measure.
\end{itemize}

%==========================================================================
\section{Conclusion}
\label{sec:conclusion}
%==========================================================================

The wealth-tax-and-growth debate is conducted as though the
question were ``does wealth tax depress growth?''  In our model the
answer depends entirely on the measure of growth, and the dependence
is first-order.  Under extractive-inclusive measures (GMT, CGDP) the
wealth tax appears to destroy 99\% of growth.  Under productive,
capability, knowledge, and median-wealth measures the wealth tax
raises growth by 9--15\%.  The pattern is robust across
productivity-channel regimes including the orthodox-dominant
calibration with capital-deepening coefficient five times the
human-capital coefficient.

The deeper finding is the CGDP growth-rate decomposition: the
$+5.76\%$ per-tick growth in status-quo CGDP is the model's
representation of the UK 2008--2024 financialisation pattern, in
which GDP grew through financial-sector compounding while productive
output, capability, and real wages stagnated.  The orthodox claim
``wealth tax depresses growth'' is mechanically correct in our model
but the growth it depresses is the measurement artefact that counted
the compounding flywheel as growth.

The implication for fiscal policy is straightforward.  The choice
between the wealth tax and the status quo is not a choice about
growth.  It is a choice about which measure of growth the
policy-maker accepts as the welfare maximand.  Under any measure
that distinguishes real production from financial-sector compounding,
the wealth tax is unambiguously growth-enhancing.  Under the
particular measure that counts financial-sector compounding as
growth, the wealth tax is growth-destroying by construction.  These
are not different empirical findings.  They are the same finding
under two different measurement choices.

We do not propose that GDP should be abandoned as a policy
instrument; it has unique uses in fiscal-capacity assessment.  We
propose that GDP should not be the welfare maximand for tax-policy
choice.  Under any alternative welfare maximand consistent with the
GDP-critique tradition the optimal-tax conclusion changes.  This is
the methodological contribution we offer.

The companion wealth-tax paper recommends a recurrent wealth tax of
5--10\% above \pounds 10M.  The present paper provides the empirical
justification for treating that recommendation as growth-favorable
under any non-extractive-inclusive welfare measure.  The
orthodox-Mirrlees prediction that the recommendation would be
growth-depressing is correct only on the measure the Mirrlees
framework was constructed around.

\bibliographystyle{plainnat}
\bibliography{references}

\appendix
\section{Productivity-channel sensitivity}
\label{app:sensitivity}

The four productivity-channel regimes tested in
Section~\ref{sec:robust} span the conventional growth-accounting
range for both the human-capital and capital-deepening coefficients.
Full numerical results are available in
\texttt{out/data/growth\_portfolio.json}.  Robustness over
$\beta_H \in [0.0, 0.7]$ and $\beta_C \in [0.0, 0.7]$ has been
tested; the headline pattern is invariant.

\section{Calibration of the tick-year mapping}
\label{app:tickyear}

The model has no intrinsic tick-year mapping.  Reported per-tick
growth rates can be annualised at any chosen ratio.  The CGDP growth
rate of $+5.76\%$ per tick under status quo, if a tick is interpreted
as one quarter, gives an annualised rate of approximately $+25\%$,
which is too high.  If a tick is interpreted as one year, the rate
is exactly $+5.76\%$, which is approximately the UK financial-sector
gross value added growth rate over 2008--2024 (though not the headline
GDP growth rate of $+1\%$).  The model's qualitative pattern (CGDP
grows fast under status quo via extractive compounding, collapses
under wealth tax) is invariant to the tick-year choice.  The
quantitative match to specific UK time-series targets requires
calibration of the tick-year ratio as a free parameter, which is left
for Path-C work.

\end{document}
